\section{Introduzione}

\subsection{Scopo del documento}
Lo scopo del documento è descrivere l'analisi dei requisiti del progetto \textbf{Second Brain}. Il testo definisce in modo strutturato i casi d’uso e le funzionalità richieste, classificandole in requisiti obbligatori, desiderabili e opzionali, secondo le specifiche fornite dal capitolato Zucchetti.

\subsection{Scopo del prodotto}
Il prodotto mira a fornire un editor Markdown intelligente che supporti l'utente non solo nella scrittura, ma anche nell'analisi critica e creativa dei propri testi tramite l'integrazione di LLM.

\subsection{Glossario}
Al fine di evitare ambiguità, i termini tecnici o con significato specifico sono seguiti dal pedice e definiti nel documento \textit{Glossario v2.1.0}.

\subsection{Riferimenti}
\subsubsection{Normativi}
\begin{itemize}
	\item Capitolato C6 - Progetto Second Brain: \\ \url{https://www.math.unipd.it/~tullio/IS-1/2025/Progetto/C6.pdf}
	\item Regolamento progetto didattico: \\ \url{https://www.math.unipd.it/~tullio/IS-1/2025/Dispense/PD1.pdf}
	\item \textit{Norme di Progetto v1.1.2}
\end{itemize}

\subsubsection{Informativi}
\begin{itemize}
	\item Slide del corso di Ingegneria del Software (T05 - Analisi dei Requisiti)
	\item Verbali esterni ed interni
\end{itemize}